\section{Prescribed-field theta data, phase descent, and the original rational tests}
\label{sec:rational-theta}

This chapter works on one original geometric realization.  It closes the
finite and arithmetic constructions between its geometric class-sum data
and the literal rational tests.  It also constructs the prescribed-field
\texttt{PinsOverL} from a precisely stated remaining geometric cover packet.
The latter packet is an input here: the existence of the higher moduli
schemes, their level maps, and their algebraic character sections is not
proved by calling the corresponding source theorem.  Similarly, the
geometric point identity and the generation of cusp forms by the original
class sums remain identified producers.  All finite transfers after those
producers are proved below, including the zero-reference-lift case.

The source is fixed at commit
\texttt{7cbf7a0c32fca1badf05931879a37c3f8a2b03b4}.  Its correctness is assumed.
The statements below distinguish a complete source endpoint from a
sufficient conditional refinement of a larger endpoint.  No new geometric
realization, coefficient field, reference phase, measure, or congruence
parameter is chosen in order to bypass an original compatibility condition.

\subsection{The prescribed coefficients and the tuple data}

% BEGIN RTG1
\paragraph{RTG1 (fixed objects and the exact remaining geometric cut).}
\label{claim:RTG1}
Fix the original $D,G,A,T,R,F,\kappa,v_0$ and
\(\mathrm{FWTie}(T,R,\kappa,v_0)\), with the original
\(\mathrm{Hk3.Core}(R,F)\), tuple off/used laws, and the original base
records $O_i$ satisfying \texttt{OMPins}.  Put
\[
 L=D.\mathrm{what}_2.\mathrm{LL},\qquad
 \sigma=D.\mathrm{what}_2.\mathrm{embLL},\qquad
 M=\sigma(L)\subset\mathbb C,\qquad K_s=T.Q.c.F.                 \tag{RTG.1}
\]
For a fixed $i$, abbreviate $V=V_c(D,i)$, $W=W_c(D,i)$,
$n=\operatorname{rank}V$, $r=\operatorname{rank}W$ and
$\mu= A.\mu_0$.  The statement of \texttt{OMPins} already says
\(O_i.c'.F=M\).  The map $j_L:L\to O_i.c'.F$ is therefore the
isomorphism satisfying $\iota_Mj_L=\sigma$; surjectivity follows by writing
any element of $M$ as $\sigma(a)$ and using injectivity of $\iota_M$.
All constructions below retain this very coefficient record $c'=O_i.c'$,
not merely an unspecified field isomorphic to it.

Here is the geometric input needed for the first construction.  For each
tuple $t$, supply the original objects
\[
 d_t:\mathrm{DeltaDat}(T.Q.Xo(i+1),c',n),\quad
 U_{2,t},\quad\mathrm{dict}_{2,t},\quad\xi_t,
                                                               \tag{RTG.2}
\]
with the exact fields of \texttt{TupleCover.CoversAt}: prescribed
$\mathrm{pd}=\mathrm{tupPrime}(R,t)$, the half ideal of RTG10,
the corresponding level $l_{2,t}$, \texttt{CovPin}, \texttt{HalfIs},
\texttt{XiIs}, the one \texttt{CXiVal} on the original
$iotaW(T,i),x_0$, all \texttt{ModelW} fields, and a positive-order
root-of-unity law for each $d_t.c_\delta$.  The coefficient field in
\texttt{CoversAt} is exactly $c'$.  RTG12 proves the full \texttt{PinsOverL}
conclusion on this input.  RTG10 constructs its numerical levels and
coefficient records, but not the higher geometric models in (RTG.2).

For the rationality transfer we retain three exact geometric assertions:
\texttt{ClsEx} on the actual tuple records used below and on their
\texttt{scaleXi} modifications; \texttt{ClassSumsSpan} for the base $O_i$;
and \texttt{PointIdentCore} with one nowhere-zero function $\lambda$ and
one group element $z_0$, chosen before the tuple records and their
\texttt{XiConst} constants.  Its hypotheses use the same original
$\mathrm{gramR},c_{0R},c_{0L}$ and the same prescribed level map.
These are genuine retained geometric producers.  Their finite consumers
are proved in RTG16--RTG20.  The base \texttt{OneMapRow0}, its original
\texttt{PhaseTie}, \texttt{LUnif}, and \texttt{EvInj} are also retained.
No conclusion here asserts existence of the common outer realization.
% END RTG1

% BEGIN RTG2
\paragraph{RTG2 (local and adelic values of the original character).}
\label{claim:RTG2}
Let $E/F$ be the original CM extension, $\chi$ a finite-order tame
Hecke character of $E$, trivial on $F$, and $M\subset\mathbb C$ a
subfield.  Suppose that, outside a finite conjugation-stable set $S$,
$\chi$ is unramified and every unramified prime value belongs to
$M^\times$.  Then every local value $\chi_P(k)$, for every finite prime
$P$ and every $k\in E^\times$, belongs to $M^\times$.
If $c$ is the original continuous adelic character with
$\mathrm{ChrEq}(c,\chi)$, every value of $c$ also belongs to $M$.
In particular the original \texttt{XiVal} follows on the same data.

\emph{Proof.}
The infinite component of a character of order dividing $a>0$ is trivial:
each factor $\mathbb C^\times$ has an $a$th root, and applying the
character to that root proves the assertion.  Fix $P\in S$ and first
$0\ne b\in\mathcal O_E$, with $v_P(b)=e$.  Chinese remainders give
$a\in\mathcal O_E$ with
\[
 a\equiv b\pmod {P^{e+1}},\qquad
 a\equiv1\pmod {Q^{e+1}}\quad(Q\in S,\ Q\ne P).               \tag{RTG.3}
\]
Thus $a\ne0$, $v_P(a)=e$, and $a/b\in1+P\mathcal O_{E,P}$.
Tameness gives $\chi_P(a)=\chi_P(b)$ and $\chi_Q(a)=1$ for
$Q\in S\setminus\{P\}$.  The principal-idele product rule now writes
$\chi_P(b)$ as the inverse of a finite product
\(\prod_{Q\notin S}\chi(Q)^{v_Q(a)}\), so it lies in $M^\times$.
Write any $k\in E^\times$ as a quotient of two nonzero integral
elements to obtain the assertion at $P$.  Outside $S$ it is the
unramified formula $\chi_P(k)=\chi(P)^{v_P(k)}$.

For the actual character $\chi=R.\xi_H(t)$, take
\[
 S=R.S_b\cup R.Q(t)\cup R.Q(\mu)
       \cup\overline{R.S_b\cup R.Q(t)\cup R.Q(\mu)}.
\]
The original ramification field gives unramifiedness off $S$.
At a prime moved by conjugation, \texttt{fam\_\allowbreak pl\_\allowbreak onto} identifies it
with a permitted source prime, and \texttt{fam\_\allowbreak muVal} gives the two
values of $\mu_{0,i}\xi_t$ and $\mu_{0,i}\xi_\mu$ in $M$.
Their quotient is $\xi_t(P)$ because the original base character
$\xi_\mu=1$.  At a conjugation-fixed prime choose
$b\in P\setminus P^2$.  The norm $b\bar b$ is integral in $F$;
triviality on $F$, and unramifiedness at the unique prime above the
underlying $F$-prime, give $\xi_t(P)^2=1$.  Such a value is $1$ or
$-1$, again in $M$.  This checks every hypothesis of the preceding
argument from the original \texttt{lev} and \texttt{fam} fields.
In particular each residue character $R.\xi_R(t,P)$ takes values in
$M$: lift its nonzero residue to an integral $b$ and apply its literal
local-value equality.  This proves the source
\texttt{xiH\_\allowbreak locK\_\allowbreak mem} and the residue-value assertion without changing
the character or using an unspecified reciprocity comparison.

For completeness, pass from local rational arguments to the original
adelic character, retaining its \texttt{ChrEq}.  The actual base-changed
adele algebra is the usual product of its infinite and finite factors.
One verifies the finite identification by choosing an $F$-basis of $E$:
at a finite place $v$ the algebra
$E\otimes_FF_v$ is the product of the completions $E_P$ above $v$.
The factors are obtained by the relatively prime local factors of a
separable defining polynomial.  Away from finitely many denominators
and its discriminant, the same map identifies the integral lattices.
It therefore gives precisely the restricted-product topology, with
the source's local-place idempotents as the individual projections.
At infinity the CM hypothesis gives a finite product of
$\mathbb C$; these identifications preserve the original rational
inclusion and the local elements used in \texttt{ChrEq}.

The finite products of ideles supported at one prime and having a
rational value $k\in E^\times$ there are dense in the finite ideles.
Indeed $E$ is dense in each completion, and a basic restricted-product
neighborhood restricts only finitely many places; enlarge this set to
contain the places where the given idele is not a unit.  Approximate
at those places and set all other components to $1$.
If $\chi^a=1$, \texttt{ChrEq} says that $c^a$ is $1$ on this
dense subgroup, hence on all finite ideles by continuity.
Rational triviality then gives $c^a=1$ on the infinite components of
rational elements.  They are dense in the infinite multiplicative
group: in a rational basis of $E$, the archimedean coordinate map is
an invertible real linear map, rational coordinates are dense, and
the set with no zero component is open.  Thus $c^a=1$ globally.
Every infinite idele has an $a$th root in the product of
$\mathbb C^\times$, so $c=1$ on the infinite factor.

The image of $c$ is now contained in the finite group of $a$th roots
of unity.  Continuity makes its kernel open.  A kernel neighborhood
in the finite ideles contains a product of open neighborhoods at
finitely many primes and all integral units elsewhere.  Given an
arbitrary finite idele $x$, enlarge that finite set to include its
nonunit places.  At each chosen prime select $k_P\in E^\times$
so close to $x_P$ that the quotient belongs to the required kernel
neighborhood.  The finite product $z$ of these supported rational
ideles has $xz^{-1}$ in the kernel.  Therefore
\[
 c(x)=c(z)=\prod_P\chi_P(k_P)\in M.
\]
This is an equality of values, not an argument taking a limit inside
the potentially nondiscrete field $M$.  It proves the full source
\texttt{char\_\allowbreak mem} implication and hence \texttt{xiVal} for $R.\xi$.
% END RTG2

% BEGIN RTG3
\paragraph{RTG3 (original character and rational-form values).}
\label{claim:RTG3}
The character of $\mathrm{levJ}(R,i,t)$ takes values in $M$ for every
$t$.  Moreover, for the original map $c_{0L,i}$ and every rational
source vector $y$, the actual definite automorphic value at any finite
adele is $c_{0L,i}(v)$ for some $v\in\mathrm{VV}(L,\lambda_i)$.

\emph{Proof.}
For $P\in R.P_g$, choose an invertible matrix with any prescribed
unit determinant, for example $\operatorname{diag}(u,1,\ldots,1)$.
The rank is positive by the original rank inequalities.  Insert that
matrix at $P$ and the identity at the other primes, and apply
\texttt{grp\_\allowbreak chiW}.  The resulting value of $R.\chi(i+1,P)$ is
$\sigma$ of the original $L$-valued finite-group character.
The level character is exactly the product of the inverses of these
block-determinant values and the direct $\xi_R(t,P)$ values from RTG2;
its unipotent extension is $1$.  Hence it is $M$-valued, with these
original distinct sign conventions.  The exponent $-1$ in the later
$\omega_Q$ residue polynomial is a separate convention.

Use the actual class presentation in $R.B.L(i,t)$, writing a finite
adele as $\gamma h_c k$.  The presentation formula for the original
rational form is
\[
 (R.HV(i,t)).e(R.jV(i,t)y)(\gamma h_c k)
   =\chi_t(k)\,R.cV_i.\rho(\gamma)\,
                       \mathrm{coef}_c(y).                 \tag{RTG.4}
\]
The original \texttt{CoefTie} supplies
$\mathrm{coef}_c(y)=c_{0L,i}(v)$.
The \texttt{CoefMapTie} in the same \texttt{FamTie} sends this, under
the rational matrix $\gamma$, to $c_{0L,i}(v')$ with the prescribed
polynomial action.  Finally write $\chi_t(k)=\sigma(a)$ and absorb
$a$ into $v'$.  Formula (RTG.4) is the original class-presentation
formula, so this conclusion concerns the actual $R.jV$ and $R.HV.e$.
It is not a newly chosen rational structure.  The coefficient maps and
their actions are exactly those reconstructed in CFB1--CFB13.
% END RTG3

% BEGIN RTG4
\paragraph{RTG4 (one coefficient extension before all tuples and classes).}
\label{claim:RTG4}
For any field embedding $j:L\to K'$ and $\sigma':K'\to\mathbb C$
with $\sigma'j=\sigma$, there is one semilinear map
\[
 c'_0:\mathrm{VV}(K',\lambda_i)\longrightarrow (R.cV_i).W
                                                               \tag{RTG.5}
\]
which satisfies the literal \texttt{CoefficientBaseChange} for
$c_{0L,i}$ and satisfies \texttt{PlaceEx} for every $t$ and every
finite family of class representatives.  In particular the conclusion
holds with $K'=O_i.c'.F=M$ and the original $j_L$.

\emph{Proof.}
Both fields have characteristic zero because they embed into $\mathbb C$.
CFB1 identifies the literal polynomial coefficient space $\mathrm{VV}$
with a kernel of finite rational linear equations in its fixed
multidegree, including the product in its separate variable blocks.
Gaussian elimination of those same equations over $\mathbb Q$, then
over $L$ and $K'$, shows that coefficientwise application of $j$
takes $\mathrm{VV}(L,\lambda_i)$ into
$\mathrm{VV}(K',\lambda_i)$.  This is also immediate on the source's
tensor generators: coefficient maps commute with each product and
determinant, and hence with every linear combination of its
dual-Weyl polynomial factors.  The argument concerns the exact
polynomial subspace, rather than an isomorphic representation.

Choose an $L$-linear projection from the full polynomial space onto
$\mathrm{VV}(L,\lambda_i)$ which is the identity on that subspace.
For each monomial $X^a$ prescribe the image
$c_{0L,i}(\pi X^a)$; extend to the $K'$ polynomial space by
$\sigma'$-semilinearity and restrict to $\mathrm{VV}(K',\lambda_i)$.
Every polynomial has finite support.  Expanding its coefficients proves
\[
 c'_0(j_*v)=c_{0L,i}(v),                                    \tag{RTG.6}
\]
including the actual signed determinant-weight space, since the
covariance and determinant factors are preserved by $j$.  In the
prescribed-field case $j_L$ is an isomorphism, so the same map is simply
$c_{0L,i}$ composed with coefficientwise $j_L^{-1}$; this description
shows that its values on the entire coefficient space are fixed.

RTG3 gives, for each $y,t,h_c$, a vector $v_c$ whose image is the
actual value at $h_c$.  Put $x_c=j_*v_c$ and use (RTG.6).
This proves \texttt{PlaceEx}.  The construction of $c'_0$ preceded
$t$, the class set, its representatives, and $y$.  It proves the whole
source \texttt{CoefExtension.exists\_\allowbreak ext\_\allowbreak placeEx} on its stated
original premises, with no extra intertwiner input.
% END RTG4

% BEGIN RTG5
\paragraph{RTG5 (the actual adelic side and a coefficient retie).}
\label{claim:RTG5}
Fix actual representatives for the rational-left, level-right double
cosets, and a coefficient map $c'_0$.  There exists a function
$\mathrm{thAd}$ satisfying the full \texttt{ThAdIs}, and its domain
predicate can be chosen to be the literal existence of a rational source
vector giving the prescribed representative values.

\emph{Proof.}
If two vectors $y,y'$ in the actual $R.HV(i,t)$ have the same values
at all representatives, their associated forms agree at
$\gamma h_c k$ by the original left and level transformation laws.
These are all finite adeles.  Injectivity of the original $R.HV.e$
therefore gives $y=y'$.  For each coefficient tuple $x$, if
$\mathrm{Place}(y,x)$ has a witness, define
\[
 \mathrm{thAd}(x)=\operatorname{onQ}
       \bigl(R.\mathrm{thF}(i,t)(y)(v_0)\bigr);
                                                               \tag{RTG.7}
\]
otherwise define it to be zero.  Uniqueness makes (RTG.7) valid for
every witness, which is precisely \texttt{ThAdIs}.  Set
$\mathrm{Plc}(x)\iff\exists y_0\,\mathrm{Place}(R.jV(y_0),x)$.
RTG4 gives \texttt{PlaceEx} on this same data.

There is also an exact retie operation on an already supplied tuple
record with the same coefficient record as the base $O_i$.  Replace only its coefficient map by (RTG.5), its level map by
the prescribed original base level map, and its \texttt{thAd/Plc} by
(RTG.7).  Preserve $c',j_L$, every geometric cover, polynomial map,
representative, weight, and scalar.  The original base
\texttt{LevIs} and level actions supply the level fields, and the
preceding argument supplies every changed representative field.
All other \texttt{Pins} fields are unchanged.  The retie thus preserves
literal coefficient-record equality, not just inclusion of fields.
In the application to RTG12 the level map is already the base map, so
this operation changes no level action at all.
% END RTG5

% BEGIN RTG6
\paragraph{RTG6 (representatives supported at the permitted split primes).}
\label{claim:RTG6}
For every tuple, every positive old bad integer $b$, and every later
level integer $M_2>0$, one can choose the actual finite representatives
in \texttt{RepsR} with $J'=\mathrm{levJ}(R,i,t)$, volume $1$, and
support at split places avoiding $\mathrm{bad}\,M_2$.  The new bad
integer is a positive multiple of $2b$ and of all the residue-prime
norms in $R.P_g\cup R.Q(t)$.

\emph{Proof.}
Take
\[
 \mathrm{bad}=2b\prod_{P\in R.P_g\cup R.Q(t)}\mathrm N(P).
                                                               \tag{RTG.8}
\]
The original self-dual frame conjugates principal congruence matrices
modulo $\mathrm{bad}$ into the level group: at a flagged prime their
reduction is the identity, so it preserves every flag and acts trivially
on the required residue quotients.  Elsewhere it preserves the original
completed lattice.  This checks the needed congruence subgroup directly.

Apply the finite factorization proved in BTH13, (BTH.49), with precision
$\mathrm{bad}^2M_2$ to each representative of the already supplied finite
class presentation $R.B.L(i,t)$.  Its rational factor is absorbed on the
left and its congruence factor on the right.  The remaining representative
is supported at finitely many split places away from
$\mathrm{bad}M_2$.  Its pure finite lift has infinite component $1$;
the unused infinite factor does not enter the definite finite class
presentation.  Distinct double cosets stay distinct, since the changes
were by exactly the two quotienting groups.  Stabilizers are conjugate
to the original finite stabilizers, so their actual cardinalities
$\gamma_c$ are positive.  Set $L_c=\mathrm{hLat}(V,h_c)$,
$\gamma_c=|\mathrm{stab}(h_c)|$, and $\mathrm{vol}=[J:J]^{-1}=1$.
These choices supply every field of \texttt{RepsR}.

For the source's literal \texttt{SplitOffBad} idempotent, select one
prime of $E$ above each supporting split place of $F$.  The product
decomposition $E\otimes_FF_v=F_v\times F_v$ gives an idempotent for
that selected factor.  The finite sum $e$ has $e\bar e=0$ and
$e+\bar e$ equal to the support idempotent.  Consequently
$(e+\bar e)(h_c-1)=h_c-1$, and each selected prime avoids
$\mathrm{bad}$.  This is the required original support identity,
including both conjugate factors.
% END RTG6

% BEGIN RTG7
\paragraph{RTG7 (the original conjugate lattices and all six representative laws).}
\label{claim:RTG7}
At an original level $M_2$, which is odd, the representatives in RTG6
admit the original
$\mathrm{La}_c=\operatorname{conjLat}(L_c)$ and integers $N_c$
satisfying all of \texttt{RepsOK} and the literal \texttt{LatPin}.

\emph{Proof.}
The completed lattice of $L_c$ is $h_c\widehat{\mathcal O_E^n}$.
This follows by intersecting with $E^n$: localization at each prime
recovers the indicated lattice by Chinese remainders, and a rational
vector integral at all primes is integral.  A common denominator embeds
this intersection between two full standard lattices, proving finite
generation and full rank.  The standard Hermitian form is self-dual;
unitarity of $h_c$ therefore makes $L_c$ self-dual.  Conjugating the
coordinates converts the source pairing convention to
$\sum\bar v_jw_j$, and gives exactly \texttt{GenLat.Unimod} for
$\operatorname{conjLat}(L_c)$.

At primes over $\mathrm{bad}M_2$ the completed lattice is standard.
The finite module measuring the difference from the standard lattice
therefore has an annihilator $d\in\mathbb N$ coprime to
$\mathrm{bad}M_2$: take a product of sufficiently high powers of the
rational primes in its support.  The same $d$ can be chosen for both
inclusions.  It supplies \texttt{TrivAt bad} and \texttt{TrivAt $M_2$}.
Because $2\mid\mathrm{bad}$, $d$ is odd; put $N_c=M_2d$.
Then $N_c$ is odd, $M_2\mid N_c$, and
\(N_c\mathcal O_E^n\subset L_c\cap\mathcal O_E^n\), also after
conjugation.  These are exactly the three \texttt{LevOK} clauses.

For \texttt{SplitSupp}, choose a sufficiently large power of the ideal
of the supporting split $F$-primes.  Its ideal class has finite order,
so a further power is principal, generated by an integral $a\ne0$.
Finiteness of the ideal class group follows from the ideal-class
representative bound already checked in KB7.  The principal ideal has
no other prime factors.  Viewed in $\mathcal O_E$, multiplication by
$a$ clears both lattice inclusions, and every prime dividing $a$ is
moved by conjugation.  This is the source's element-valued
\texttt{SplitSupp}, rather than only a statement about local support.
The positive stabilizer cardinality was proved in RTG6.  Finally the
source \texttt{LatPin} says exactly that membership in $\mathrm{La}_c$
is membership of the conjugate vector in $L_c$; this is its definition.
Thus unimodularity, both triviality clauses, level, positive cardinality,
and split support all hold on the original lattices.
% END RTG7

% BEGIN RTG8
\paragraph{RTG8 (the exact square law by paired local elementary divisors).}
\label{claim:RTG8}
Let $\mathcal L\subset E^n$ be self-dual for the standard Hermitian
pairing and have the source \texttt{SplitSupp}.  Suppose an integral
ideal $\mathfrak c$ satisfies
\[
 \mathfrak c+\bar{\mathfrak c}=\mathcal O_E,\qquad
 \det(\mathcal L)\bar{\mathfrak c}=\mathfrak c.
                                                               \tag{RTG.9}
\]
Then, for some integer $k\ge1$,
\[
 [\mathcal L:\mathcal L\cap\mathcal O_E^n]
       =\mathrm N(\mathfrak c)k^2.                            \tag{RTG.10}
\]
If $r=2s+1$ is odd, the positive rational number
\[
 q=\bigl(\mathrm N(\mathfrak c)^s k^{2s+1}\bigr)^{-1}
 \quad\hbox{satisfies}\quad
 q^2[\mathcal L:\mathcal L\cap\mathcal O_E^n]^r
       =\mathrm N(\mathfrak c).                             \tag{RTG.11}
\]

\emph{Proof.}
At a prime $P$ write the elementary divisors relative to the standard
lattice as integers $a_{P,1},\ldots,a_{P,n}$.  This is ordinary row and
column elimination over the discrete valuation ring: choose an entry
of smallest valuation, clear its row and column, and induct.  Applied
to a common denominator of a basis matrix it also allows negative
integers.  The determinant valuation is $d_P=\sum_j a_{P,j}$,
and the length of
$\mathcal L_P/(\mathcal L_P\cap\mathcal O_{E,P}^n)$ is
$e_P=\sum_j\max(0,-a_{P,j})$.
Self-duality says that a matrix for the lattice at $\bar P$ is the
conjugate inverse transpose of a matrix at $P$, up to an integral
invertible basis change.  Thus its elementary divisors are the negatives
of those at $P$, in reverse order.  A prime fixed by conjugation does
not meet the supporting element in \texttt{SplitSupp}; both clearing
inclusions then give $\mathcal L_P=\mathcal O_{E,P}^n$.

For a conjugate pair $\{P,\bar P\}$, put $q_P=\mathrm N(P)=\mathrm
N(\bar P)$.  Its contribution to the index is
$q_P^{\sum_j|a_{P,j}|}$.  Equation (RTG.9) and coprimality say that its
contribution to $\mathrm N(\mathfrak c)$ is
$q_P^{|\sum_j a_{P,j}|}$.  The difference of these exponents is a
nonnegative even integer: if $A$ and $B$ are the sums of the positive
and of the absolute negative entries, it equals
$A+B-|A-B|=2\min(A,B)$.  Multiplying over the finitely many supporting
pairs proves (RTG.10).  The index-product formula itself follows by a
composition series of the finite torsion quotient; a factor
$\mathcal O_E/P$ contributes its cardinality $\mathrm N(P)$.
The ideals in (RTG.9) are nonzero, so all quantities used in (RTG.11)
are positive.  Direct substitution gives (RTG.11).  If $n=0$, the
same argument gives $\mathfrak c=\mathcal O_E$, index $1$, and $k=q=1$.
This proves the complete source \texttt{law\_\allowbreak of\_\allowbreak squares}, with its
actual index and its arbitrary odd exponent.
% END RTG8

% BEGIN RTG9
\paragraph{RTG9 (the ideal-valued coefficient and the original \texttt{EpIs}).}
\label{claim:RTG9}
Starting from the base \texttt{BIs}, every tuple has a function
$B_t:\operatorname{Ideals}(\mathcal O_E)\to M$ with the exact
\texttt{BIs}; and there is one function $\mathrm{epF}_t$ of the lattice
satisfying all of \texttt{EpIs}.

\emph{Proof.}
For a prime away from the new bad integer, the original base identity
and $\xi_\mu=1$ put
$\mu_{0,i}(P)\eta(P)^{-1}\sqrt{\mathrm N(P)}$ in $M$.
Multiply by the actual value $\xi_t(P)\in M$ from RTG2.  Unique
factorization of a nonzero integral ideal gives
\[
 B_t(\mathfrak c)=
  (\mu_{0,i}\xi_t\eta^{-1})(\mathfrak c)
                   \sqrt{\mathrm N(\mathfrak c)}\in M
                                                               \tag{RTG.12}
\]
when $\mathfrak c$ is coprime to the new bad integer; multiplicativity
of the positive real square root verifies the formula prime by prime.
Set $B_t=0$ off that guard.  This proves precisely \texttt{BIs}, with
no assertion at excluded ideals.

The ideal $\mathfrak c$ in (RTG.9), if it exists for a given lattice,
is unique: at each conjugate pair its exponents are
$\max(d_P,0)$ and $\max(-d_P,0)$, and at a fixed prime they vanish.
For a lattice satisfying the actual three-part \texttt{Guard} and
having such $\mathfrak c$, take $q$ from RTG8 and put
\[
 \mathrm{epF}_t(\mathcal L)=B_t(\mathfrak c)^{-1}q.
                                                               \tag{RTG.13}
\]
Off this existential guard set it to zero.  Uniqueness of $\mathfrak c$
makes the same value valid for every later admissible determinant ideal.
RTG8 gives the required positive rational square identity, and
(RTG.13) is the exact remaining equality in \texttt{EpIs}.
The source exponent is $T.Q.Dw(i+1).p+T.Q.Dw(i+1).q=r$;
the original signature and rank laws say it is odd.  Finally set
$\mathrm{ep}(t,c)=\mathrm{epF}_t(L_c)$, keeping the original
unconjugated $L_c$ in this particular definition.
% END RTG9

% BEGIN RTG10
\paragraph{RTG10 (the exact half levels and their coefficient fields).}
\label{claim:RTG10}
Under the original \texttt{QGuard}, let $I_t$ be the product of the
extensions to $E$ of the distinct $K$-primes under $R.Q(t)$, and let
$d_t^{\rm num}$ be the product of its distinct rational residue primes.
Then the source half level exists with
\[
 M_{2,t}=M_0d_t^{\rm num},\quad
 M_{b,2,t}=M_{b,0}I_t,\quad A_{2,t}=1,
 \qquad I_t\bar I_t=(d_t^{\rm num}),\quad I_t+\bar I_t=1.
                                                               \tag{RTG.14}
\]
Moreover its coefficient record can be chosen on exactly $M$, with
its specified root of unity, and no enlargement of the base coefficient
field.

\emph{Proof.}
A residue prime in $R.Q(t)$ has residue degree one and is unramified
in the original $E$-datum.  It avoids $2$.  In the quadratic field
$K=\mathbb Q(\sqrt{-d})$ its prime therefore has degree and ramification
index one, so it is one of the two distinct conjugate primes.  More
explicitly, reducing the original generator $w_0$ gives an integer $e$
with $e^2-e+(d+1)/4\equiv0\pmod q$ and $1-2e\not\equiv0\pmod q$.
The two ideals are $(q,w_0-e)$ and $(q,w_0-(1-e))$; their sum is $1$
and their product is $(q)$.  Extension to $E$ preserves these two
identities.  The first clause of \texttt{QGuard} says that two selected
primes above the same rational $q$ have the same $K$-prime below them.
Thus the distinct $K$-primes and the distinct rational primes index
the same factors, rather than counting a rational prime several times.
This proves the two ideal identities in (RTG.14).

Each $q$ is odd and different from the two base primes $g_0,g_1$.
The original base level has $M_0=g_0g_1$ and $A_0=1$.
Consequently $d_t^{\rm num}$ is odd and coprime to $M_0$.
The new half ideal is generated over $\mathcal O_E$ by elements in
$\mathbb Z[w_0]$, as is the old one.  Products preserve this property.
Together with (RTG.14) this checks every field of the literal
\texttt{Level} and \texttt{IsCoverLevel}: oddness, $M_{2,t}\ge3$,
the generator condition, product, coprimality, and the three divisibility
conditions for the old level.  The old half ideal is the product of
$R.P_g$: it is a squarefree divisor of $(g_0g_1)$ with exactly those
prime factors, by the original half pins.  Hence \texttt{HalfIs} holds,
including $I_t\subset P$ for $P\in R.Q(t)$ and the two bad-prime
avoidance clauses.

The original \texttt{fam\_\allowbreak roots} supplies a primitive $q$th root in
$L$ for each new rational prime.  Products of primitive roots of
coprime orders have product order: if $(xy)^a=1$, raise to each of the
two orders and use coprimality.  Induction produces a primitive
$d_t^{\rm num}$th root and, together with the existing primitive
$M_0$th root, a primitive $M_{2,t}$th root in $M$.
Every complex $M_{2,t}$th root is a power of this one, since the
polynomial $X^{M_{2,t}}-1$ has exactly these distinct roots.
Thus the particular root prescribed by the coefficient record belongs
to $M$.  Use the original inclusions and this root to form $c_{2,t}$
with $c_{2,t}.F=M$; its field comparison with $c'$ is the identity.
This constructs coefficients and levels only.  It does not construct
$X_{o,2,t}$, the level morphisms, or their algebraic sections.
% END RTG10

% BEGIN RTG11
\paragraph{RTG11 (the exact residue polynomial and its quotient descent).}
\label{claim:RTG11}
For any ideal $J$ with $\bar J\subset P$ for every
$P\in R.P_g\cup R.Q(t)$, there is an $M$-valued function $\varphi_t$
on the original lattice quotient satisfying the full \texttt{PhiIs},
with the source exponent $-1$ on $\omega_Q$.

\emph{Proof.}
Use the support idempotent supplied by the original \texttt{thP.vP},
and keep the actual lattice embedding $\iota=iotaW(T,i)$.
For $x=(x_b)$ in that lattice the matrix is literally
\[
 U(x)=e\,\beta_V^{-1}
       \bigl(\operatorname{gram}(V)^{-1}\bar X
                          \operatorname{gram}(W)\bigr)\beta_W,
 \qquad X_{ba}=\iota(x_b)_a.                                \tag{RTG.15}
\]
The standard-space condition gives $\operatorname{gram}(V)=1$.
The original frame has denominator $1^c=1$, so its rational lattice
is the original $\mathrm{hLat}(V,1)=\mathcal O_E^n$.
A self-dual frame of the standard completed lattice is integral with
integral inverse: both matrices carry a basis of the same lattice to
another such basis.  For the source lattice one must also retain its
scalar factor: $\operatorname{range}(\iota)=\mathrm{latG}(W,1)$ is the
$\mathcal O_E$-span of $a v$ with
$a\in AA=\mathfrak c_E\mathfrak b^{-1}$ and $v$ in the original
unscaled completed lattice.  Put $N=\mathrm N(\mathfrak b)>0$.
The containment $N\in\mathfrak b$ follows because multiplication
by $N$ kills the finite group $\mathcal O_E/\mathfrak b$.
Thus $Na\in\mathfrak c_E\subset\mathcal O_E$.
The original host law says that every rational prime dividing $N$
divides $Q_Sq$.  At a selected prime, \texttt{HalfIs} makes
$Q_Sq\,\mathrm N(\mathfrak c_K)$ a unit, and therefore makes $N$
a unit.  It follows that $\bar a$ is integral there.
Pairing $v$ with a self-dual $\beta_W$ column is integral there;
multiplication by $\bar a$ preserves integrality.  At all other
places the support idempotent is zero.  Additivity and conjugate
$\mathcal O_E$-linearity extend the argument across the entire
span defining $\mathrm{latG}(W,1)$.  Together with the integral
$\beta_V^{-1}$ this proves integrality of (RTG.15), including
the exact inverse and conjugation positions.

If $x_b$ is changed by an element of $J\iota(\mathrm{La})$, its
conjugate changes by $\bar J\bar\iota(\mathrm{La})$.
The same integral pairing calculation therefore changes every selected
residue matrix by zero.  Thus
\[
 \omega_K\bigl(R.P_g,\overline{U(x)}\bigr)
 \ \omega_Q\bigl(R.Q(t),\operatorname{rank}W,
                      R.\xi_R(t),-1,\overline{U(x)}\bigr)
                                                               \tag{RTG.16}
\]
is independent of the quotient representatives.  Here the overline
on $U(x)$ in (RTG.16) denotes its family of residue matrices, not a
second complex conjugation.  The first factor is $0$ or $1$; the
second is either zero or a finite product of inverse character values.
RTG2 therefore places (RTG.16) in $M$.  Define $\varphi_t$ to be that
value.  Integrality and (RTG.16) are exactly the two conclusions in
\texttt{PhiIs}; no averaging normalization is inserted.
For the application $J=\mathrm{Mfr}(l_{2,t})=\bar M_{b,2,t}$,
(RTG.14) gives the required inclusion at every selected prime.
This is the complete source \texttt{exists\_\allowbreak phi\_\allowbreak core} on its stated
original hypotheses.
% END RTG11

% BEGIN RTG12
\paragraph{RTG12 (complete \texttt{PinsOverL} on the exact cover packet).}
\label{claim:RTG12}
On the original premises of RTG1 and the geometric \texttt{CoversAt}
packet (RTG.2), there is an all-tuple record $\widetilde O_i$ with
\[
 \widetilde O_i.c'=O_i.c',\qquad
 \widetilde O_i.j_L=O_i.j_L,\qquad
 \widetilde O_i.\mathrm{lev}=O_i.\mathrm{lev},\qquad
 \widetilde O_i.c'.F=M,                                    \tag{RTG.17}
\]
and every field of the original \texttt{Pins}.  Hence the literal
\(\mathrm{PinsOverL}(T,R,\kappa,v_0,i)\) holds.

\emph{Proof.}
Select once the original \texttt{FamTie} witness
$\mathrm{gramR},c_{0R},c_{0L}$ and use RTG4 for one $c'_0$.
For each tuple, RTG6--RTG7 give the representatives and all the exact
lattices at its already selected $d_t.l_2$; RTG9 gives $B_t$ and
$\mathrm{epF}_t$; RTG11 gives $\varphi_t$; and RTG5 gives
$\mathrm{thAd}_t,\mathrm{Plc}_t$ on these representatives.
Keep $O_i.\mathrm{io}$ and $O_i.\mathrm{hw}$ unchanged.  Define the
source weight by the literal finite transform
\[
 w_{t,c}(q)=e_{F,t}^{-1}
   \left(\widehat{\,e_{F,t}\varphi_t\,}(q)\right).             \tag{RTG.18}
\]
It is independent of $c$.  The defining \texttt{WPin} equality follows
by applying $e_{F,t}$ to (RTG.18); no Fourier inversion theorem is
needed for this field.

For completeness the eleven \texttt{Pins} fields are as follows.
The \texttt{lift} field uses the original field/standard-space/vector/
level identities and RTG4.  The \texttt{tuple} field uses the
representatives, \texttt{PlaceEx}, \texttt{ThAdIs}, the exact
\texttt{Plc} equivalence, $B_t$, the supplied \texttt{XiIs}, RTG11,
and the supplied half identity.  Both \texttt{act} fields are the
unchanged base actions on the same coefficient record and level map.
The \texttt{bad} field follows because the new integer is a multiple
of the original bad integer.  The \texttt{reps} field is RTG7, and
\texttt{ep} is RTG9 with its defining equality on each class lattice.
The \texttt{io} field is the original base field.  The \texttt{w} field
is (RTG.18).  Finally \texttt{cov}, \texttt{cxi}, and \texttt{model$_2$}
are exactly the supplied geometric cover fields, on the same
$iotaW,x_0$ and the same coefficient inclusion.  This list accounts for
every field, including all tuple-dependent ones.

RTG12 is a sufficient conditional reconstruction of
\texttt{pinsOverL\_\allowbreak of\_\allowbreak clauses}.  It does not count the implication
\texttt{CoverField}$\Rightarrow$\texttt{CoversAt} as proved: that
implication still needs actual higher models, level maps, and character
sections.  The numerical levels, coefficient records, representatives,
coefficient maps, residue functions, scalar normalizations, and adelic
side have all been constructed independently here.
% END RTG12

\subsection{Class sums and the prescribed phase}

% BEGIN RTG13
\paragraph{RTG13 (the original evaluation map, including its fallback case).}
\label{claim:RTG13}
An original \texttt{FWTie} together with $R.\mathrm{LUnif}$ satisfies
the full \texttt{FWTie.Evaluates}.

\emph{Proof.}
Fix $i$, an embedding $\tau$ with the required \texttt{GoodEmb}, and
the literal equality identifying the target space.  If a semilinear
evaluation map exists, the \texttt{ev} field of \texttt{FWTie} gives
the claimed equality for its original $j$.  Otherwise its \texttt{ev0}
field says that this original $j$ is zero.  The independence clause of
\texttt{LUnif} then forces its rational coordinate space to be zero:
any nonzero vector would be a linearly independent singleton whose image
is zero.  The isomorphism $\mathrm{fw}.e$ consequently forces every
original cusp vector to be zero.  The zero semilinear map is then an
evaluation map, because evaluating the zero section and the zero
multiplicity map both gives zero.  This contradicts the fallback case.
Thus the original $j$ evaluates correctly at every permitted embedding.
The proof checks the zero-dimensional case rather than assuming a
nonzero cusp form.  It is the complete source
\texttt{evaluates\_\allowbreak of\_\allowbreak unif} argument on its original hypotheses.
% END RTG13

% BEGIN RTG14
\paragraph{RTG14 (independence and scalar descent on the actual cusp reads).}
\label{claim:RTG14}
Assume the original \texttt{LUnif} and \texttt{EvInj}.  Let
$C_i=T.\mathrm{cuspA}(i+1,\mathrm{wtA}(\kappa_i),\kappa_i.n)$ over
$K_s$, and let $\psi(f)$ be its original adelic read, with the original
component factor $T.c_E$ but without the positive factor
$\ell_i=\mathrm{lieFac}>0$.  A $K_s$-independent finite family in $C_i$
has $\mathbb C$-independent images under $\psi$.
Consequently, if two nonzero finite $M$-linear combinations of these
reads differ by a scalar $a\in\mathbb C$, then $a\in M$.

\emph{Proof.}
The original $\mathrm{fw}.e$ identifies $C_i$ with the exact
$K_s$-form $(R.FW_i).WK$.  In \texttt{IsKsForm}, a $K_s$-basis of
that form is an $L$-basis of the surrounding rational coordinate space.
To see that any independent family remains independent over $L$,
extend it to a $K_s$-basis; its change-of-basis matrix has nonzero
determinant in $K_s$, hence still has nonzero determinant in $L$.
\texttt{LUnif} takes this family to a complex-independent family under
the original $j$.  By RTG13 its evaluation at the original $v_0$ is
$\ell_i\psi(f)$.  The original \texttt{EvInj} makes this evaluation
injective, and $\ell_i\ne0$, proving the first assertion.

For the second, choose a $K_s$-basis $b_1,\ldots,b_d$ of $C_i$;
finite dimensionality is an original \texttt{LawsJ} field.
Rewrite the two combinations as $\sum u_j\psi(b_j)$ and
$\sum v_j\psi(b_j)$, with all $u_j,v_j\in M$, since the image of
$K_s$ lies in $M$.  Independence gives $u_j=av_j$ for each $j$.
A nonzero first combination has some $u_j\ne0$, so $v_j\ne0$ and
$a=u_j/v_j\in M$.  This is the complete finite \texttt{ratio\_\allowbreak mem}
argument, now installed on the original cusp reads; no unproved general
adelic linear-independence theorem has been used.
% END RTG14

% BEGIN RTG15
\paragraph{RTG15 (the read of a class sum uses the unchanged base comparison).}
\label{claim:RTG15}
Suppose that a tuple record has exactly the base $c'$ and level map,
as in (RTG.17), and the base \texttt{OneMapOf} holds with the prescribed
component function and phase $p$.  For every cusp form $f$ over $M$,
its actual read is a finite $M$-linear combination $\Psi_f$ of the
base reads in RTG14 and, for every $z\in G$, satisfies
\[
 c_Q(q)\,\mathrm{rd}_z(f)(q)
     =\chi_G(z)\,\Psi_f(q^{\rm out}k_z).                     \tag{RTG.19}
\]
Here $k_z$ is any original level lift of $z$, and
\[
 \begin{split}
 c_Q(q)&=T.c_E(i+1,\sigma_0)(\mathrm{dict}(q^{\rm out})),\\
 \chi_G(z)&=\prod_{P\in R.P_g}\chi_{i+1,P}(\det R.e_G(z)_P).
 \end{split}
\]
All these factors and maps are the original ones.

\emph{Proof.}
The base \texttt{OneMapOf} already contains the full \texttt{CuspBC}
and \texttt{ActBC} for this very coefficient record and level map.
Thus write $f$ as a finite base change $\sum a_j f_j$, where
$a_j\in M$ and $f_j\in C_i$.  Apply \texttt{ActBC} and the definition
of evaluating a base-changed section at a frame.  The base
\texttt{EquivAt} says, term by term, that reading
$\mathrm{actLev}(z)f_j$ equals $\chi_G(z)$ times reading $f_j$
at the right-translated point.  Summing gives (RTG.19) with
$\Psi_f=\sum\iota_M(a_j)\psi(f_j)$.
This argument uses exactly the \texttt{IsBC} equality on frames; its
finite sum commutes with the original reading and level action.
It does not assume a new base-change theorem for a different field.

These functions are rational-left-invariant.  Indeed the original
multiplicity-map range is a span of automorphic maps, each of which has
that invariance; RTG13 identifies its evaluation with
$\ell_i\psi(f_j)$.  Finite combinations preserve the property.
An equality of two such functions at all $q^{\rm out}k$ is therefore
an equality on the whole adelic group: for any $g$, let $q$ represent
$gk^{-1}$, so $g=\gamma q^{\rm out}k$ for a rational $\gamma$.
The factor $\chi_G(z)$ never vanishes, since its factors are units.
% END RTG15

% BEGIN RTG16
\paragraph{RTG16 (the exact finite class-sum proof of \texttt{LiftsSpan}).}
\label{claim:RTG16}
Use the base $O=O_i$ and an all-tuple $O'$ with (RTG.17).
Assume \texttt{ClassSumsSpan} for $O$, \texttt{ClsEx} for $O'$, and
\texttt{PtIdent} for both records with the same $\lambda,z_0$ and
nonzero constants $c_0,c_1$.  Then the literal
\(\mathrm{LiftsSpan}(R,i,v_0)\) holds.

\emph{Proof.}
Write $\theta_t(y)=\operatorname{onQ}(R.\mathrm{thF}(i,t)
(R.jV(i,t)y)(v_0))$.  \texttt{PlaceEx} and \texttt{ThAdIs} identify
this with $O'.\mathrm{thAd}_t(x)$ for an actual coefficient tuple
$x$ satisfying \texttt{Plc}.  \texttt{ClsEx} gives a cusp form $f_t$
with the exact \texttt{IsThClsD} value identity.  Because the coefficient
record is literally the same, $f_t$ belongs to the very $M$-space
occurring in \texttt{ClassSumsSpan}.  Thus, by the definition of linear
span, there are finitely many $a_j\in M$ and base class sums $f_j$
with $f_t=\sum a_j f_j$.  Each $f_j$ comes from an actual rational
base source vector $y_j$ by the \texttt{Plc} equivalence.

Let $s_t$ and $s_0$ be the root-of-unity factors in the two point
identities.  They are fixed before $x$ or $f_t$, and $s_0\ne0$.
The identities and linearity of the unchanged $\mathrm{rd}_{z_0}$ give
\[
 \begin{split}
 c_1\theta_t(y)&=s_t\lambda\,\mathrm{rd}_{z_0}(f_t),\\
 c_0\theta_\mu(y_j)&=s_0\lambda\,\mathrm{rd}_{z_0}(f_j),\\
 \theta_t(y)&=
   \frac{s_tc_0}{c_1s_0}\sum_j\iota_M(a_j)\theta_\mu(y_j).
 \end{split}                                                  \tag{RTG.20}
\]
The last equality follows by substituting the second line into the
first and dividing by $c_1s_0$; it does not divide by a theta value.
This proves membership in the complex span of the original base lifts.
It also proves that all tuple lifts vanish if all base lifts vanish.
This is the exact class-sum-to-\texttt{LiftsSpan} implication on the
prescribed-field route, with no dimension comparison and no added rapid
decay hypothesis.
% END RTG16

% BEGIN RTG17
\paragraph{RTG17 (the phase ratio lies in the prescribed field).}
\label{claim:RTG17}
Under RTG14--RTG16, suppose some base lift is nonzero.  If the common
point identity uses constants $c_0$ on $O$ and $c_1$ on $O'$, then
\(c_0/c_1\in M^\times\).  Multiplying only the original cover section
$d_t.c_\delta$ by $e_{F,t}(c_0/c_1)$ gives an all-tuple record with
all its \texttt{Pins} intact and \texttt{XiConst} equal to $c_0$.

\emph{Proof.}
Take a rational base source vector $y$ with $\theta_\mu(y)\ne0$.
Represent it in both records using their \texttt{PlaceEx} and
\texttt{ThAdIs}.  The base \texttt{OneMapOf} gives a class-sum cusp
form $f_0$; the all-tuple \texttt{ClsEx} gives $f_1$.  Their point
identities have factors $s_0,s_1\in M^\times$ and say
\[
 c_0\theta_\mu(y)=s_0\lambda\,\mathrm{rd}_{z_0}(f_0),\qquad
 c_1\theta_\mu(y)=s_1\lambda\,\mathrm{rd}_{z_0}(f_1).
                                                               \tag{RTG.21}
\]
Cancel the nowhere-zero $\lambda$.  Then use (RTG.19), cancel the
nonzero $\chi_G(z_0)$ and the component factor, and extend the
result from $q^{\rm out}k_{z_0}$ by rational-left-invariance.  The
result is
\[
 \Psi_{f_0}=\frac{c_0s_1}{c_1s_0}\Psi_{f_1}.                 \tag{RTG.22}
\]
The first function is nonzero, by (RTG.21).  RTG14 therefore puts the
ratio in (RTG.22) in $M$, and multiplication by $s_0/s_1$ proves the
claim.  This step is where the original coefficient-field identity
matters.

Set $u=c_0/c_1\in M^\times$ and perform exactly \texttt{scaleXi}:
$d_t.c_\delta\mapsto e_{F,t}(u)d_t.c_\delta$, with all other data
fixed.  The level-section space is linear, and evaluation multiplies
by the same nonzero scalar, so regularity and nonvanishing survive.
All the other clauses of $d_t.\mathrm{IsLev}$ are unchanged.
The \texttt{XiConst} equality is multiplied by $u$ and hence now has
constant $uc_1=c_0$.  Both constants are roots of unity by the original
\texttt{CXiVal}; the new constant therefore has precisely the required
root-of-unity law.  The covariance and \texttt{cxi} fields have just
been checked; the other nine \texttt{Pins} fields are copied unchanged.
In particular no coefficient map, representative, theta
map, weight, lattice, frame, or level action has been rescaled.
% END RTG17

% BEGIN RTG18
\paragraph{RTG18 (complete all-tuple \texttt{OneMapOf} with the original phase).}
\label{claim:RTG18}
On the geometric cut of RTG1, the strengthened \texttt{PinsOverL}
of RTG12 and the base \texttt{PhaseTie}/\texttt{OneMapRow0} give
\texttt{OneMapOverL} with the original prescribed phase $p$.
This statement preserves every \texttt{OneMapOf} field.

\emph{Proof.}
Apply \texttt{PhaseTie} to the supplied base \texttt{OneMapRow0}.
This gives the base \texttt{OneMapOf} with the actual $T.c_E$ and
$p$, rather than a newly selected phase.  Choose the coefficient
witness from the base \texttt{lift.c0} first.  RTG4--RTG5 retie the
all-tuple record to it, keeping (RTG.17).  Apply the one retained
\texttt{PointIdentCore} to the two records; it supplies the same
$\lambda,z_0$ before their constants $c_0,c_1$ are selected from
\texttt{CXiVal}.  The retained \texttt{ClsEx} applies to both the
resulting tuple record and, if needed, its \texttt{scaleXi} modification.

If there is a nonzero base lift, RTG17 gives that modification and the
new point identity has $c_1=c_0$.  Keep it as $O'$.  If all base lifts
are zero, keep the unmodified $O'$.  In either case let $z_1$ be the
central group element, and $s_*$ the root of unity, in the base
\texttt{OneMapRow}.  For a tuple class sum $f_t=\sum a_j f_j$ as in
RTG16, its base equations are
\[
 p\theta_\mu(y_j)=
   s_*\,c_Q\ell_i\,\mathrm{rd}_{z_1}(f_j).
                                                               \tag{RTG.23}
\]
If $c_1=c_0$, multiply (RTG.20) by $p$ and substitute (RTG.23):
\[
 p\theta_t(y)=
   \left(s_*\frac{s_t}{s_0}\right)
       c_Q\ell_i\,\mathrm{rd}_{z_1}(f_t).                    \tag{RTG.24}
\]
Its tuple scalar lies in $M$ and is a root of unity; a common multiple
of the three positive orders proves this directly.  It is selected
before $x$ and the class sum.  If every base lift is zero, RTG16 makes
$\theta_t(y)=0$, while (RTG.23) and the same finite expansion make
the right side of (RTG.24) zero.  Thus (RTG.24) remains valid, with the
same root scalar, without asserting $c_0/c_1\in M$ in that case.

\texttt{ClsEx} supplies an actual cusp form for each permitted $x$ and
says that every section with the same \texttt{IsThClsD} values equals
it.  Hence (RTG.24) holds for every such section, as required by
\texttt{OneMapRow}, not just for a chosen representative.  The component
norms, $\lVert p\rVert=1$, and centrality of $z_1$ are copied from the
base.  The coefficient inclusion, \texttt{CuspBC}, and \texttt{ActBC}
are literally the base fields because $c'$ and the level map are
unchanged.  \texttt{EquivAt} is the same original action/read identity.
Finally \texttt{PhRat} refers only to the base tuple, the unchanged
$\mathrm{thG}$, $v_0$, $c_Q$, and $p$, so it is copied from the base
\texttt{OneMapOf}.  This verifies the complete record and proves the
asserted \texttt{OneMapOverL}.

This is a conditional proof of the selected
\texttt{oneMapOverL\_\allowbreak of\_\allowbreak clauses} route: the exact common point identity,
class-sum existence/uniqueness, and base class-sum span are the geometric
suppliers retained in RTG1.  They have not been replaced by a generic
isomorphism of spaces.
% END RTG18

\subsection{Rational theta coordinates and the actual lattice}

% BEGIN RTG19
\paragraph{RTG19 (\texttt{OneMapOverL} implies the actual \texttt{RatAt}).}
\label{claim:RTG19}
For the original \texttt{FWTie}, \texttt{IsModels}, \texttt{IsThP},
\texttt{LUnif}, and \texttt{EvInj}, assume
\(\mathrm{OneMapOverL}(T,R,\kappa,v_0,p,i)\).
Then for every $t$ and every original rational source vector $x$
there is $y\in L^{R.nW_i}$ with
\[
 \begin{split}
 (R.HW_i).e\bigl((R.FW_i).j(\sigma)y\bigr)
      &=p\,R.\mathrm{thF}_{i,t}(R.jV_{i,t}x),\\
 (R.FW_i).j(\sigma)\bigl((R.FW_i).e_{\rm op}y\bigr)
      &=p\,R.\mathrm{thG}_{i,t}(R.jV_{i,t}x).
 \end{split}                                                  \tag{RTG.25}
\]
These are the full source \texttt{RatAt} clauses.

\emph{Proof.}
Choose the all-tuple $O$ in \texttt{OneMapOverL}.
\texttt{PlaceEx} gives the actual coefficient tuple of $R.jVx$;
\texttt{ThAdIs} identifies its adelic theta value; and
\texttt{OneMapRow} gives its cusp class sum and a scalar $\zeta\in M$.
The \texttt{CuspBC} field expands that class sum into finitely many
base cusp forms.  The \texttt{ActBC} and \texttt{EquivAt} fields move
the single prescribed central action through the expansion.  Thus the
actual equality, at every quotient representative, is
\[
 p\,R.\mathrm{thF}_{i,t}(R.jVx)(v_0)(q^{\rm out})
    =\sum_j\zeta a_j\,\ell_i\,
             O.\phi(T.c_E)(f_j)(q^{\rm out}),\qquad a_j\in M.
                                                               \tag{RTG.26}
\]
This is the finite calculation of RTG15 with the scalar
$\zeta$ and the positive $\ell_i$ restored.  Every $f_j$ is a cusp
form over $K_s$, including after the level action.
Since $M=\sigma(L)$, choose $b_j\in L$ with
$\sigma(b_j)=\zeta a_j$ and set
\[
 y=\sum_j b_j\,\mathrm{fw}.e(f_j)\in L^{R.nW_i}.              \tag{RTG.27}
\]
RTG13 gives the exact evaluation of $\mathrm{fw}.e(f_j)$.
Here $\mathrm{sg0}(D,A)=\sigma$, its restriction is the original
$s_0$, and the field equality identifying $O.\mathrm{s0}$ is the
original coefficient inclusion.  Therefore no conjugation of $b_j$
or additional Lie factor occurs in (RTG.26).

The \texttt{thF\_\allowbreak mem} field puts the actual theta map in the original
$R.HW.e$ range; write it as $R.HW.e(w)$.  All elements of this range
are rational-left-invariant after evaluation, since it is the specified
span of automorphic multiplicity maps.  Equation (RTG.26) therefore
holds on the full adelic group by the quotient argument of RTG15.
\texttt{EvInj} now implies $R.FW.j(y)=p w$, proving the first line
of (RTG.25).  For the second line one must retain the original good projection.
Write $e_g$ for the coefficients of $\mathrm{idemOf}(A.X)$.
The literal operator is
\[
 e_{\rm op}=\sum_{g\in G}e_g\chi(g)^{-1}\rho(g).
 \tag{RTG.27a}
\]
The original \texttt{RhoIs} in \texttt{LUnif} says that
$R.HW.e\circ R.FW.j\circ\rho(g)$ equals $\sigma(\chi(g))$
times right action by $g$ on $R.HW.e\circ R.FW.j$.
Multiplying by $\sigma(e_g\chi(g)^{-1})$ cancels the character
and summing gives
$R.HW.e(R.FW.j(e_{\rm op}y))=e_W(R.HW.e(R.FW.j(y)))$,
where $e_W$ is the original finite group-algebra projector.
The printed \texttt{thP.thG} field is
$R.HW.e(\mathrm{thG}(x))=e_W(\mathrm{thF}(x))$.
Combining it with the first line of (RTG.25) and injectivity of
$R.HW.e$ proves the second line.  Thus the projected rational vector
is $e_{\rm op}y$; the proof does not identify it with the unprojected
$y$, or replace the original projector by an adjoint.
The proof is exactly on the hypotheses of
\texttt{ratAt\_\allowbreak of\_\allowbreak oneMap}/\texttt{rat\_\allowbreak of\_\allowbreak oneMapOverL}; RTG13
supplies its printed evaluation hypothesis whenever \texttt{LUnif}
is given.
% END RTG19

% BEGIN RTG20
\paragraph{RTG20 (the phase-correct \texttt{RatG} family).}
\label{claim:RTG20}
Apply RTG12 and RTG18 for every original target index $i$, on the
same original data and the geometric cut in RTG1.  Then RTG19 gives
\[
 \forall i\ \forall t\ \forall x\quad\exists y\in L^{R.nW_i}:
 \quad R.FW_i.j(\mathrm{sg0})y
          =p_i\,R.\mathrm{thG}_{i,t}(R.jV_{i,t}x).
                                                               \tag{RTG.28}
\]
This is the literal \texttt{RatG}.  In the three-stage family take
\[
 p_i=\operatorname{cj}\bigl(\mathrm{twAt}(\sigma,A_3.A.v_0)\bigr)
                         (R.\mathrm{phs}_i),                 \tag{RTG.29}
\]
as required by the original \texttt{PhaseTie}.

\emph{Proof.}
RTG12 chooses one coefficient map before all tuples, representatives,
and vectors.  The common point identity in RTG1 chooses
$\lambda_i,z_{0,i}$ before the tuple records and constants.
RTG17, when applicable, makes only one scalar correction on each
all-tuple record, simultaneously for all its tuples.  RTG18 retains
$p_i$ from the base phase tie.  Finally RTG19 is universal in $t,x$;
use its projected rational vector $e_{\rm op}y$ in (RTG.28).
Only this output vector depends on $t,x$.  These are exactly
the quantifiers in (RTG.28), and substituting (RTG.29) yields the
original family proposition.  No coefficient field or phase is selected
after $x$.  No common threshold or universal family existence is
inferred from this pointwise construction.
% END RTG20

% BEGIN RTG21
\paragraph{RTG21 (the original separating tests at both Boolean levels).}
\label{claim:RTG21}
For every original \(R:\mathrm{RealDat3}(D,A_3,K_s)\), face $F$,
and $b\in\{\mathrm{false},\mathrm{true}\}$,
\texttt{LUnif} and the \texttt{RatG} with phase (RTG.29) imply
\[
 \mathrm{LiftTestB}(R.\mathrm{toCx3}(F),b),
 \qquad \mathrm{LiftTest}(R.\mathrm{toCx3}(F)).                \tag{RTG.30}
\]

\emph{Proof.}
Put $s=\mathrm{twAt}(\sigma,A_3.A.v_0)$.
The original $\mathrm{LUnif}$ says that the vectors
$R.FW_i.j(\mathrm{sg0})e_k$ are independent and span the actual
$R.HW_i.H$.  Twisting them by $\mathrm{Tw.mk}_s$ gives a complex
basis $b_k$ of the actual space $\mathrm{Tw}_s(R.HW_i.H)$.
For $y\in L^{R.nW_i}$ the original embedding identity is
\[
 \sigma(a)=\operatorname{cj}_s\bigl(A_3.A.v_0.\mathrm{embedding}(a)\bigr).
\]
By the twisted scalar action and $\operatorname{cj}_s^2=1$,
\[
 \mathrm{Tw.mk}_s(R.FW_i.j(\mathrm{sg0})y)
     =\sum_k A_3.A.v_0.\mathrm{embedding}(y_k)\,b_k.
                                                               \tag{RTG.31}
\]
Choose the coordinate functionals of this one basis.  They separate
points, including when its dimension is zero.  The actual
$X.\theta'$ in $X=R.\mathrm{toCx3}(F)$ is the twist of
$p_iR.\mathrm{thG}$ with precisely (RTG.29).  Equations
(RTG.28)--(RTG.31) show that every coordinate of
$X.\theta'(X.j_Lx)$ belongs to the range of
$A_3.A.v_0.\mathrm{embedding}$.

For the extra level-$b$ assertion, take the actual rational basis
$R.bL(F,b,i)$ of the original good subspace.  The definition of $X.w0C$
gives
\[
 \mathrm{Tw.out}_s\bigl(X.w0C(b,i)e_\ell\bigr)
       =R.FW_i.j(\mathrm{sg0})\bigl(R.bL(F,b,i)_\ell\bigr).
                                                               \tag{RTG.32}
\]
Indeed the embedding takes the coordinate vector $e_\ell$ to itself,
and the defining linear combination reduces to its $\ell$th basis
vector.  By (RTG.31), its $k$th separating coordinate is
\[
 A_3.A.v_0.\mathrm{embedding}(R.bL(F,b,i)_{\ell k}).
\]
This proves the last \texttt{LiftTestB} field at the original level.
The same basis of the ambient target was selected before $t,x$ and
before this level basis.  Dropping the last field gives
\texttt{LiftTest}.  Thus the complete source
\texttt{liftTestB\_\allowbreak of\_\allowbreak ratG} is reconstructed for both Boolean values;
the base-false conclusion used by the main route is a specialization.
% END RTG21

% BEGIN RTG22
\paragraph{RTG22 (the full original \texttt{MK} from the actual valuation lattice).}
\label{claim:RTG22}
On a given original \texttt{FWTie}, the original \texttt{LatFull}
field implies $(R.FW_i).\mathrm{MKFull}$ for every $i$.

\emph{Proof.}
The actual norm and point correspondences in \texttt{FWTie} give,
for $f\in C_i$,
\[
 \mathrm{fw}.e(f)\in(R.FW_i).MK
  \quad\Longleftrightarrow\quad
   v(f)\ge0\text{ for every original finite valuation and point}.
                                                               \tag{RTG.33}
\]
For clarity, the forward implication evaluates the norm at the inverse
image of any original point under \texttt{fw.pt}; the norm identity
and \(\|f\|_v\le1\iff v(f)\ge0\) give the displayed inequality.
The reverse implication uses the forward point correspondence in the
same way.  This proves (RTG.33) on exactly the original measures of
integrality, not only on a smaller chosen set of points.

\texttt{LatFull} says that this valuation-nonnegative lattice in $C_i$
is finitely generated over $\mathcal O_{K_s}$ and spans $C_i$ over
$K_s$.  The isomorphism $\mathrm{fw}.e:C_i\simeq(R.FW_i).WK$ and
(RTG.33) identify its image with the whole original $MK$.
The image of a finite generating set generates, and the image of its
span is $WK$.  These are exactly the two \texttt{MKFull} fields.
This proof uses the original \texttt{LatFull}; in the common selected
realization it is the named \texttt{InB.latFull}/geometric fullness
input.  It is not deduced here from \texttt{FWTie} alone.
% END RTG22

% BEGIN RTG23
\paragraph{RTG23 (the literal coordinate lattice at every prescribed level).}
\label{claim:RTG23}
For any original $R,F$ with \texttt{IsKsForm} and \texttt{MKFull},
and every Boolean $b$, the literal lattice predicate
\(\mathrm{MLat}(R.\mathrm{toCx}(F),b)\) holds.
Together with RTG22 this proves the source
\texttt{GoodLat.mLat\_\allowbreak of\_\allowbreak tie} on its printed \texttt{LatFull} premise.

\emph{Proof.}
First put $N=\mathcal O_L\,MK$ in the actual rational coordinate space
$L^{R.nW_i}$.  A finite $\mathcal O_{K_s}$-generating set of $MK$
also generates $N$ over $\mathcal O_L$, since the inclusion
$\mathcal O_{K_s}\to\mathcal O_L$ is the original one.
The \texttt{IsKsForm} basis and the $K_s$-span statement for $MK$
show that $LN$ is the full coordinate space.  Thus $N$ is a full
$\mathcal O_L$-lattice.

Here is the exact intersection argument, including finite generation.
If $N\subset V$ is a full lattice and $f:U\hookrightarrow V$ is an
injective $L$-linear map, then $f^{-1}N$ is a full lattice in $U$.
Its image is a submodule of the finite module $N$ and hence finite
over the noetherian ring $\mathcal O_L$; injectivity transfers a finite
generating set to the preimage.  For $u\in U$, expand $f(u)$ in a
finite $L$-linear combination of elements of $N$ and clear the finitely
many denominators.  Some nonzero $a\in\mathcal O_L$ has
$af(u)\in N$, so $au\in f^{-1}N$ and
$u=a^{-1}(au)$ lies in its $L$-span.

Apply this to the original coordinate map
\[
 f_b:L^{d_b}\longrightarrow L^{R.nW_i},\qquad
          c\longmapsto\sum_\ell c_\ell R.bL(F,b,i)_\ell.
                                                               \tag{RTG.34}
\]
It is injective because $R.bL$ is the specified basis of the original
$W_{\rm good}(J_b,\lambda_b,e_C)$.  Its image already satisfies the
level and good-subspace conditions.  Therefore the source's
\texttt{Mcoord} is exactly $f_b^{-1}N$, including both those conditions;
no saturation or replacement lattice is introduced.  The preceding
argument proves its finite generation and full span.  By the definition
of $R.\mathrm{toCx}$ this is its original $M(b,i)$, so it is the
required \texttt{MLat}.  Zero-dimensional good subspaces are included.
No archimedean metric is changed by this rational lattice argument.
% END RTG23

% BEGIN RTG24
\paragraph{RTG24 (the base tests and lattice on one original three-stage family).}
\label{claim:RTG24}
On the same $R:\mathrm{RealDat3}$, $F$, $T$, $\kappa,v_0$, and
base $O_i$, assume the original algebraic/model/phase data of RTG1,
the exact retained geometric cut, \texttt{LUnif}, \texttt{EvInj}, and
\texttt{LatFull}.  Then
\[
 \begin{split}
 &\mathrm{LiftTest}(R.\mathrm{toCx3}F),\\
 &\mathrm{LiftTestB}(R.\mathrm{toCx3}F,\mathrm{false}),\qquad
   (R.\mathrm{toCx3}F).\mathrm{MLat}(\mathrm{false})
 \end{split}                                                  \tag{RTG.35}
\]
hold, with the analogous two predicates at \texttt{true} as well.

\emph{Proof.}
RTG20 provides exactly the phase-correct \texttt{RatG} consumed by
RTG21.  RTG22--RTG23 provide the coordinate lattice on the same
underlying $R.\mathrm{toRealDat}$ and face $F.\mathrm{lv2}$.
The definitions of \texttt{toCx3} preserve these particular maps,
bases, lattices, and the embedding at $A_3.A.v_0$.  Applying the proved
statements, rather than making new choices, gives (RTG.35).
All finite and all-tuple choices have already been made before these
tests.  This is the complete consumer conclusion of
\texttt{liftTests\_\allowbreak and\_\allowbreak mLat\_\allowbreak of\_\allowbreak ties} on the explicit sufficient cut
printed here.  It is not a claim that its larger original
\texttt{TiesB/TiesG} premises already have full human proofs.
% END RTG24

% BEGIN RTG25
\paragraph{RTG25 (what has closed and the smallest geometric producers still visible).}
\label{claim:RTG25}
The preceding implications have the following exact audit scope.
RTG2--RTG5 close local and adelic character-value descent, all-tuple coefficient
extension, the actual adelic side, and the selected same-coefficient
retie.  RTG6--RTG11 close split representatives, all their lattice laws,
the square normalization, ideal coefficients, half levels and coefficient
records, and the exact residue function.  RTG12 closes all eleven
\texttt{Pins} fields conditional on the specified \texttt{CoversAt}.
RTG13--RTG21 close the finite evaluation, class-span, scalar-ratio,
phase, \texttt{OneMapOverL}, \texttt{RatG}, and both-level test
implications.  RTG22--RTG24 close the full original lattice and test
installation on the named \texttt{LatFull} input.

There remain the following actual producers on the same $T,R,F$.
First, \texttt{coversAt\_\allowbreak over\_\allowbreak field} must produce the higher moduli
and level/character-section packet on $O_i.c'.F=M$; the original
\texttt{CoverField} only supplies algebraic closedness for an already
existing model.  The source route passes through
\texttt{moduliScheme\_\allowbreak levelMaps\_\allowbreak heckePoints'},
\texttt{LevelCover.exists\_\allowbreak deltaDat\_\allowbreak isLev}, and
\texttt{XiFun.exists\_\allowbreak sec\_\allowbreak char\_\allowbreak val}.

Second, the actual class sums need their cusp membership and their
uniqueness among sections: \texttt{clsEx\_\allowbreak of\_\allowbreak core} ultimately uses
\texttt{ClsVal.\allowbreak exists\_\allowbreak cusp\_\allowbreak isThClsD\_\allowbreak unique}, with the actual
\texttt{PsiTie}, polynomial and higher-cover dictionary.

Third, \texttt{point\_\allowbreak identity\_\allowbreak core} needs the common geometric
normalization and the theta/short-series identity, before its
$\lambda,z_0$ can be used for all records and scalar modifications.

Fourth, \texttt{classSumsSpan\_\allowbreak core}, with its original size inequality
$8fr+4\le n_V$, must prove generation of the original base cusp space;
a generation statement for an abstract
isomorphic space would not suffice.  These are retained substantial
geometric reconstruction tasks, not possible counterexamples to the
accepted formal proof.

The actual initial geometry, base \texttt{OMPins},
\texttt{OneMapRow0}/phase compatibility, common \texttt{FWTie}, and
\texttt{CoverField} remain where they occur in the outer common
construction.  If \texttt{LatFull} is not already supplied as an original
geometric field, its independent producer remains
\texttt{SupFull.fullAt\_\allowbreak of\_\allowbreak pins}: CM-point detection and bounded
denominators, followed by generation by short theta series and the
power-clearing argument.  RTG22 does not claim to prove these geometric
facts from \texttt{FWTie} alone.  The actual arithmetic eta agreement
needed for full \texttt{ReadL} remains the same RDR22 guard; this chapter
does not identify it with a later freely selected eta.  The stronger
second \texttt{DivRows3} theta square, primitive rows, and the common
selector duties also remain separate.  Thus no whole common-selection
or \texttt{BASE\_\allowbreak THETA\_\allowbreak RATIONAL\_\allowbreak TESTS} producer is declared closed.
What is closed is the explicit finite and arithmetic route from these
literal geometric suppliers to (RTG.35).
% END RTG25
